\providecommand{\bC}{\mathbb{C}}
\providecommand{\bZ}{\mathbb{Z}}
\providecommand{\bQ}{\mathbb{Q}}
\providecommand{\Aut}{\mathrm{Aut}}
\providecommand{\Auts}{\mathrm{Aut}_s}
\providecommand{\Fix}{\mathrm{Fix}}
\providecommand{\DT}{\mathrm{DT}}
\providecommand{\red}{\mathrm{red}}
\providecommand{\BPS}{\mathrm{BPS}}
\providecommand{\BKM}{\mathrm{BKM}}
\providecommand{\cat}{\mathrm{cat}}
\providecommand{\ch}{\mathrm{ch}}
\providecommand{\fiber}{\mathrm{fiber}}
\providecommand{\CoHA}{\mathrm{CoHA}}
\providecommand{\Borch}{\mathrm{Borch}}
\providecommand{\wt}{\mathrm{wt}}
\providecommand{\sdim}{\mathrm{sdim}}
\providecommand{\gr}{\mathrm{gr}}
\providecommand{\cO}{\mathcal{O}}
\providecommand{\cM}{\mathcal{M}}
\providecommand{\cY}{\mathcal{Y}}
\providecommand{\cBPS}{\mathcal{BPS}}
\providecommand{\fg}{\mathfrak{g}}
\providecommand{\fn}{\mathfrak{n}}
\providecommand{\kch}{\kappa_{\mathrm{ch}}}
\providecommand{\kBKM}{\kappa_{\mathrm{BKM}}}
\providecommand{\kcat}{\kappa_{\mathrm{cat}}}
\providecommand{\kfib}{\kappa_{\mathrm{fiber}}}

\section*{Order-two CHL positive half and the BPS denominator}
\label{sec:N2-CHL-positive-half-BPS}

\paragraph{The free CHL quotient.}
Let $S$ be a projective K3 surface and let
$g_2\in\Auts(S)$ be a symplectic involution.  Nikulin's theorem gives
$\Fix(g_2)$ as an eight-point set, and Mukai identifies the cohomological
class with the Mathieu class $2A$.  Let $E$ be an elliptic curve and let
$t_2\in E[2]\setminus\{0\}$.  The cyclic group
$\Gamma_2=\langle(g_2,t_2)\rangle$ acts on $S\times E$ by
\[
  (s,e)\longmapsto (g_2s,e+t_2),
\]
and the action is free because translation by $t_2$ acts freely on $E$.
The order-two CHL threefold is
\[
  X_2=(S\times E)/\Gamma_2 .
\]
The quotient map $\pi:S\times E\to X_2$ is finite \'{e}tale of degree
$2$, and the product form $\omega_S\wedge\omega_E$ descends because
$g_2^*\omega_S=\omega_S$ and $t_2^*\omega_E=\omega_E$.  Thus $X_2$ is
Calabi-Yau in the trivial-canonical sense used here; the elliptic base
contributes $h^{0,1}(X_2)=1$.

\paragraph{Etale K3 fibration over the elliptic base.}
Let $E'=E/\langle t_2\rangle$ and let
\[
  p_E:X_2\longrightarrow E'
\]
be the induced projection.  For $e'\in E'$ with inverse image
$\{e,e+t_2\}$, the fibre is
\[
  p_E^{-1}(e')
  =
  (S\times\{e,e+t_2\})/\Gamma_2
  \cong S .
\]
Thus $p_E$ is an isotrivial K3 fibration split by the \'{e}tale cover
$E\to E'$.  The singular surface quotient $S/\langle g_2\rangle$ and
its K3 resolution belong to the projection toward $S/\langle g_2\rangle$
and to the fixed-locus twining calculation.  The fibres of $p_E$ are
the smooth surfaces $S$.

\paragraph{Hodge supertrace and the four invariants.}
The finite \'{e}tale quotient gives
\[
  H^q(X_2,\cO_{X_2})
  =
  H^q(S\times E,\cO_{S\times E})^{\Gamma_2}.
\]
Since $g_2$ is symplectic on $H^{0,2}(S)$ and $t_2$ is a translation on
$H^{0,1}(E)$, each of
$H^{0,0}$, $H^{0,1}$, $H^{0,2}$, and $H^{0,3}$ contributes one invariant
line.  Hence
\[
  \kcat(X_2)=\chi(\cO_{X_2})=1-1+1-1=0 .
\]
The compact CY$_3$ Hodge-supertrace reading of the Vol III
CY-to-chiral functor gives the same number for the chiral output:
\[
  \kch(\Phi_3(X_2))=\sum_{q=0}^3(-1)^qh^{0,q}(X_2)=0 .
\]
The dimension is $d=3$, so $\Phi_3(X_2)$ is an $E_1$-chiral output.
Braiding belongs to the Drinfeld centre or to a completed
Hall-Drinfeld double after the centre, negative-half, Cartan, and
Hopf-pairing data have been supplied.

\paragraph{The order-two Borcherds coefficient.}
The primitive CHL Borcherds input is the $2A$-twined K3 elliptic genus
in the normalization $Z^{(2A)}_{K3}=2\phi^{(2A)}_{0,1}$.  Its leading
Fourier term is
\[
  Z^{(2A)}_{K3}(\tau,z)\big|_{q^0}
  =
  \zeta^{-1}+8+\zeta .
\]
Hence
\[
  c_2(0)=8 .
\]
The value is pinned three ways.  First, the same integer is Nikulin's
fixed-locus Euler number $\chi(S^{g_2})=8$ for the symplectic
involution.  Second, the Vol III Gritsenko-Nikulin ladder records
\[
  (c_1(0),c_2(0),c_3(0),c_4(0),c_6(0))
  =(10,8,6,4,2).
\]
Third, the executable table in
\texttt{compute/lib/k3\_yangian\_borcherds\_weight\_theta\_refinement.py}
has, at $N=2$,
\[
  T_{H^*}(g_2)=16,\qquad A_2=4,\qquad
  c_2(0)=T_{H^*}(g_2)-2A_2=8 .
\]
The singly twined $2A$ Jacobi convention records central coefficient
$4$ before the factor $Z^{(2A)}_{K3}=2\phi^{(2A)}_{0,1}$ is restored.
The primitive CHL denominator uses the restored coefficient $8$.
Therefore
\[
  \kBKM(\Delta^{(2)}_{4,\BKM})
  =
  \wt(\Delta^{(2)}_{4,\BKM})
  =
  \frac{c_2(0)}{2}
  =
  4 .
\]
The fixed-locus number is also the order-two fibre witness
\[
  \kfib^{(2)}=\chi(S^{g_2})=8 .
\]
Thus the four order-two quantities are
\[
\begin{array}{c|c|c}
\hbox{invariant} & \hbox{value} & \hbox{source}\\ \hline
\kcat(X_2) & 0 & \hbox{finite \'{e}tale descent and K\"unneth}\\
\kch(\Phi_3(X_2)) & 0 & H^{0,\bullet}(X_2)\ \hbox{supertrace}\\
\kfib^{(2)} & 8 & \chi(S^{g_2})\\
\kBKM(\Delta^{(2)}_{4,\BKM}) & 4 & c_2(0)/2
\end{array}
\]
These numbers live on distinct constructions.

\paragraph{Untwisted $K3\times E$ constants.}
The unquotiented $K3\times E$ anchor supplies a different list:
\[
  \kcat(K3\times E)=0,\qquad
  \kch(\Phi_3(K3\times E))=0,\qquad
  \kch^{\mathrm{Heis}}(K3\times E)=3,
\]
\[
  \kBKM(\Delta_5)=5,\qquad
  \kfib(K3)=24 .
\]
The first two values are compact Hodge supertraces on the threefold, the
third is the Heisenberg-Mukai chiral specialization, the fourth is the
weight of the Gritsenko $\Delta_5$ BKM denominator, and the fifth is the
Mukai-rank fibre witness.

\paragraph{Separation from adjacent automorphic rows.}
The symbol $\Delta^{(2)}_{4,\BKM}$ denotes the primitive CHL order-two
Borcherds denominator of weight $4$.  The Gritsenko-Cl\'ery
diagonal-divisor table is a separate eight-row classification: its row
$F_2=\Delta_2$ has weight $2$, and its congruence row
$Q_1$ has weight $1$ on the $(t,N)=(2,2)$ cover.  A comparison between
these automorphic products requires a row map specifying the Jacobi
input, product lattice, modular group or cover, character, Weyl chamber,
and scalar square.  The order-two CHL denominator comparison here uses
the primitive coefficient $c_2(0)=8$ and the weight-four product
$\Delta^{(2)}_{4,\BKM}$.

The Oberdieck-Pandharipande scalar belongs to the scalar reduced-DT
partition function.  It supplies a scalar automorphic normalization;
Hall multiplication, positive-half BPS data, and BKM recognition require
the additional structures below.

\paragraph{Reduced DT moduli and the positive half.}
Bryan-Oberdieck present the CHL quotient by a symplectic K3 automorphism
and an elliptic translation, and formulate the primitive reduced DT
partition functions for order two classes.  Curve classes on $X_2$ are
recorded by the descended K3 numerical class together with the elliptic
degree; integrally one keeps the Bryan-Oberdieck distinction between
untwisted and twisted primitive classes.  Let $C_2^+\subset N_1(X_2)$ be
the semigroup generated by effective such classes with nonnegative
elliptic degree.

For $\gamma\in C_2^+$, the reduced DT stack
$\cM_{\DT}^{\red}(X_2;\gamma)$ is presented on the cover by
$\Gamma_2$-equivariant reduced DT objects on $S\times E$ of class
$\pi^!\gamma$.  The reduced obstruction theory descends because the
Kool-Thomas cosection from the K3 holomorphic two-form is
$\Gamma_2$-equivariant.  Choose the resulting orientation data and write
$\cBPS^{\red}_{\gamma}$ for the oriented d-critical BPS sheaf.  The
effective Hall positive half is
\[
  \cY^+_{\BPS}(X_2)
  =
  \bigoplus_{\gamma\in C_2^+}
  H^\bullet_c\!\left(
    \cM_{\DT}^{\red}(X_2;\gamma),
    \cBPS^{\red}_{\gamma}
  \right),
\]
with product given by extension convolution in the effective cone.
This is an $E_1$ Hall algebra.  Its primitive BPS Lie positive half is
\[
  \fn^{(2),+}_{\BPS}
  =
  \mathrm{Ind}_{\BPS}\bigl(\cY^+_{\BPS}(X_2)\bigr),
\]
the BPS indecomposable object with Hall commutator bracket on the
associated graded.  This positive half contains effective charges.  The
negative half, Cartan, Hopf pairing, Weyl vector, and denominator
character are separate data.

\paragraph{Denominator comparison.}
Let
\[
  J^{\mathrm{CHL}}_2(\tau,z)
  =
  \sum_{n,\ell}c_2(4n-\ell^2,\ell)q^n\zeta^\ell
\]
be the primitive order-two CHL Jacobi input in the above normalization.
Let $L_2^+\subset \bZ^3$ be the positive charge monoid in the
Gritsenko-Nikulin order-two product chamber.  The target-side
Borcherds product is the weight-four automorphic form
\[
  \Delta^{(2)}_{4,\BKM}(Z)
  =
  e^{2\pi i(\rho_2,Z)}
  \prod_{\alpha=(n,\ell,m)\in L_2^+}
  \left(1-e^{2\pi i(\alpha,Z)}\right)^{
    c_2(4nm-\ell^2,\ell)
  },
\]
where $Z$ is the genus-two period variable and $\rho_2$ is the Weyl
vector in this chamber.  Borcherds' weight formula gives
$\wt(\Delta^{(2)}_{4,\BKM})=c_2(0)/2=4$.

A Hall-Borcherds recognition datum for $X_2$ consists of the oriented
Hall source above, a compatible negative half and Cartan, a completed
Hopf pairing, a radical quotient, parity-refined PBW, and a comparison
map from primitive BPS generators to the Borcherds charge lattice.  If
this datum exists and satisfies the denominator compatibility, the
comparison statement is
\[
  \gr\,\fn^{(2),+}_{\BPS}
  \cong
  \fn_+\!\left(\fg_{\Delta^{(2)}_{4,\BKM}}\right),
  \qquad
  \sdim\,\gr_\alpha\fn^{(2),+}_{\BPS}
  =
  c_2(4nm-\ell^2,\ell),
\]
for $\alpha=(n,\ell,m)\in L_2^+$.  The right side is the signed
superdimension convention for Borcherds root multiplicities.  The full
BKM algebra
$\fg_{\Delta^{(2)}_{4,\BKM}}$ is obtained after adjoining the negative
half, Cartan, completion, and denominator character.  The positive half
supplies the source-side primitive Lie algebra and the coefficient
comparison above.

\paragraph{Source anchors.}
Nikulin, \emph{Finite automorphism groups of K\"ahler K3 surfaces}
(1980), and Mukai, \emph{Finite groups of automorphisms of K3 surfaces
and the Mathieu group} (1988), give the order-two symplectic fixed
locus.  Bryan-Oberdieck, \emph{CHL Calabi-Yau threefolds: curve
counting, Mathieu moonshine and Siegel modular forms} (2018), give the
free CHL quotient and the order-two reduced DT scalar partition
function.  Borcherds' singular-theta weight formula and the
Gritsenko-Nikulin order-two denominator ladder give
$\kBKM(\Delta^{(2)}_{4,\BKM})=c_2(0)/2=4$.  The Gritsenko-Cl\'ery
diagonal-divisor catalogue supplies separate automorphic rows; row
identification requires the comparison datum listed above.
